\documentclass[11pt,a4paper]{article}
\usepackage[margin=25mm,headheight=15pt]{geometry}
\usepackage[T1]{fontenc}
\usepackage[utf8]{inputenc}
\usepackage{lmodern,microtype}
\usepackage{amsmath,amssymb,amsthm,mathtools,mathrsfs}
\usepackage{booktabs,array,longtable}
\usepackage{enumitem,xurl,needspace}
\usepackage[hidelinks]{hyperref}
\usepackage{fancyhdr}
\hypersetup{pdftitle={Reversion Beyond Archimedean Valuations},pdfauthor={Research article prepared for Vladimir Reshetnikov},pdfsubject={Hahn power-logarithmic reversion, support-depth Newton iteration and finite certificates}}
\pagestyle{fancy}\fancyhf{}
\fancyhead[L]{\small Reversion beyond Archimedean valuations}
\fancyhead[R]{\small Research article}
\fancyfoot[C]{\thepage}
\setlength{\parskip}{0.25em}
\setlength{\parindent}{1.2em}
\setlist{itemsep=0.25em,topsep=0.4em}
\emergencystretch=2em
\allowdisplaybreaks
\newtheorem{theorem}{Theorem}[section]
\newtheorem{proposition}[theorem]{Proposition}
\newtheorem{lemma}[theorem]{Lemma}
\newtheorem{corollary}[theorem]{Corollary}
\theoremstyle{definition}
\newtheorem{definition}[theorem]{Definition}
\newtheorem{example}[theorem]{Example}
\newtheorem{question}[theorem]{Research question}
\theoremstyle{remark}
\newtheorem{remark}[theorem]{Remark}
% ed. (2026-09-29): unnumbered environment for ProveIt editorial notes; it
% leaves every numbered statement of the delivered article unchanged.
\newtheorem*{ednote}{Editorial note (ProveIt, 2026-09-29)}
\newcommand{\N}{\mathbb N}
\newcommand{\Z}{\mathbb Z}
\newcommand{\Q}{\mathbb Q}
\newcommand{\R}{\mathbb R}
\newcommand{\HH}{\mathscr H}
\newcommand{\AS}{\mathscr A_S}
\newcommand{\supp}{\operatorname{supp}}
\newcommand{\Id}{\operatorname{Id}}
\newcommand{\ord}{\operatorname{ord}}
\newcommand{\Res}{\operatorname{Res}}
\newcommand{\depth}{\mathsf H}
\newcommand{\Div}{\mathsf D}
\newcommand{\coef}[1]{\left[#1\right]}
\newcommand{\Ocal}{\mathcal O}
\newcommand{\rising}[2]{\prod_{j=#1}^{#2}}
\newcommand{\pathcode}[1]{\nolinkurl{#1}}
\title{\textbf{Reversion Beyond Archimedean Valuations}\\[0.65em]
\large Strong Hahn convergence, sharp Newton depth,\\
and minimal finite coefficient certificates}
\author{Research article prepared for Vladimir Reshetnikov\\
\small Based on the ProveIt transseries research program}
\date{29 September 2026}
\begin{document}
\maketitle
\begin{abstract}
A positive valuation gain need not become cofinal in a non-Archimedean
ordered group. Consequently, an argument that proves reversion by sending
that gain to infinity can fail even when the inverse exists and every one
of its coefficients has a finite exact computation. We resolve the
fixed-shift Hahn power--logarithmic version of the ProveIt research question
on reversion beyond Archimedean exponent groups.

The replacement invariant is the maximum length of a positive support word
representing a prescribed exponent. The all-word form of Neumann's lemma
makes this length finite without any Archimedean hypothesis. It yields a
complete support filtration, a finite-dimensional implicit-system theorem,
and the ordinary integer-indexed Lagrange--B\"urmann inverse. For Newton
iteration we prove the sharper support-depth law $r\mapsto2r+1$ and the
uniform bound $r_n=2^{n+1}-1$, with an example attaining every bound.
We characterize exactly when valuation convergence holds uniformly for
this class: the multiples of its least support exponent must be cofinal.
In contrast, strong Hahn convergence always holds under the stated support
hypotheses.

For any finite set of requested coefficients we construct a finite-rank
quotient algebra that is minimal among algebra quotients retaining those
coefficient maps. This gives terminating exact certificates even when a
valuation cutoff contains infinitely many monomials. A rank-two infinite
support has depth $n^2+k+1$ at $(n,k)$ and admits no positive real-valued
additive grading. More generally, every positive-integer superadditive
profile is realized by such a rank-two support, allowing arbitrarily fast
depth growth. Proofs, exact symbolic checks, a statement-level formalization
plan, and further research questions are supplied. Generalized-series
implicit functions are classical; the contribution here is the explicit
repository resolution, sharp precision theory, and certificate construction,
not a claim of first discovery of Hahn-series inversion.
\end{abstract}
\newpage
\begingroup
\small
\setlength{\parskip}{0pt}
\tableofcontents
\endgroup
\newpage

\section{Research target, contribution, and scope}
\label{sec:scope}

\subsection{The repository question}
The inspected repository snapshot is ProveIt commit
\pathcode{3f2d57344dc6f0b3dfa376f9490bfe13d1beb81e}, retrieved on
29 September 2026. Its transseries development has moved to
\pathcode{Analysis/Transseries}. The recent article
\emph{Support-Controlled Reversion and the Exact One-Exponential
Substitution Group} asks, in its further-research section, what replaces
the real-valued support bound and the sequential fixed-point argument when
the exponent group is arbitrary. It explicitly separates failure of
cofinal valuation improvement from failure of existence
\cite{RepoReversion,RepoReadme}.

That is the principal target here. The same article also proposes
proof-producing support and precision inference. Our finite quotient
construction addresses a precise algebraic part of that second target:
which coefficients must be retained, why the retained set is finite, and
what finite precision suffices for exact Newton inversion.

The inspection was targeted: the recent article, the transseries README,
and the well-based-support Lean module were read. The much larger canonical
volume was not exhaustively audited. The immutable source addresses in the
bibliography and the accompanying provenance file identify the exact
materials used. We make no claim to have compared every statement against
every document in the repository.

% ed. (2026-09-29): scope note on "We resolve" and "sharper", with the
% repository antecedents the article could not see or did not cite.
\begin{ednote}\label{ed:scope}
\emph{Scope.} ``We resolve'' in the abstract means: this article answers the research
question ``Beyond Archimedean exponent groups'' of
\path{Analysis/Transseries/docs/series-and-transseries/Support_Controlled_Reversion_One_Exponential/reversion_and_one_exponential.tex}
for fixed-shift Hahn power--logarithmic derivations, as stated in this
section. It is not the first non-Archimedean Hahn
reversion in the repository: Lemma \texttt{e:lem-reversion} of the
surcomplex analysis volume \cite{EdSurcomplex}
(\path{Algebra/SurrealNumbers/docs/surcomplex/analysis/article.tex})
already proves existence, uniqueness and the same positive-integer
Lagrange sum for Hahn sections over a non-Archimedean exponent group, by
the finite-word form of Neumann's lemma, without logarithmic blocks,
characters or depth bounds; and the all-word lemma is checked in Lean
(see the note after Lemma~\ref{lem:words}).
\emph{Priority of the Newton rate.} The law $r\mapsto2r+1$,
$r_n=2^{n+1}-1$, called ``sharper'' in the abstract, is for $\Gamma\subseteq\R$ and
$h=1$ exactly the rate $e_{k+1}\ge2e_k+\alpha+2$,
$e_k\ge(2^{k+1}-1)\delta-1$ of Theorem~4.1 (\texttt{thm:newton},
equation \texttt{eq:newtonrate}) of that reversion article, expressed in
units of $\delta$; the article does not cite it. New here are its
validity without cofinality, its attainment (Theorem~\ref{thm:sharp}),
the valuation and cut criteria, the minimal quotients and the rank-two
depth profiles.
\emph{Uncited canonical remarks.} Remark
\texttt{plt:rmk:ext-archimedean} of the canonical volume
(\path{Analysis/Transseries/docs/series-and-transseries/Transseries_And_Inversion/transseries_and_inversion.tex})
shows the phenomenon of Section~\ref{sec:sharp} for lexicographic
$\Z^2$ and $S=\{(0,1)\}$, and Remark \texttt{plt:rmk:ext-neumann-general}
there quotes, without proof, the all-word clause of Neumann's lemma for
non-Archimedean groups that Lemma~\ref{lem:words} proves.
\end{ednote}

\subsection{What is classical and what is established here}
Higman's word theorem and the associated positive-support summability
principle are classical \cite{Higman}. General implicit-function methods
for generalized series already exist: van der Hoeven develops Noetherian
operators and an implicit-function theorem, including summable recursive
expansions \cite{vdH}. Composition and different modes of transseries
convergence are treated by Edgar \cite{Edgar}; a recent general Taylor
expansion framework is developed by Bagayoko and Mantova \cite{BM}.
The coefficient Lagrange inversion identity is classical as well
\cite{Gessel}.

Accordingly, the claim here is not that non-Archimedean formal inversion was
previously unknown. We give a self-contained, quantitatively precise answer
in the coefficient algebra used by the repository, with the extra data
needed to define its derivative outside a real exponent group. The central
outputs are as follows.

\begin{center}
\begin{tabular}{@{}p{0.26\textwidth}p{0.67\textwidth}@{}}
\toprule
Output & Precise content\\
\midrule
Existence and reversion & Arbitrary ordered-group supports are allowed;
ordinary integer Lagrange indices suffice; the normalized inverse stays in
the positive monoid generated by the input support.\\[0.45em]
Sharp precision & Newton error depth improves from $r$ to $2r+1$;
$r_n=2^{n+1}-1$ is attained by an explicit scalar inverse.\\[0.45em]
Topology boundary & Cofinality of multiples of the minimum support is
necessary and sufficient for valuation convergence uniformly over the
small-coefficient implicit systems considered here.\\[0.45em]
Finite certification & Each finite coefficient request factors through a
minimal finite-rank quotient; a finite-support instance has a rational
separating-functional certificate.\\[0.45em]
Infinite-support obstruction & A well-ordered rank-two support has exact
quadratic depth and no positive real additive grading.\\
\bottomrule
\end{tabular}
\end{center}

All these statements are proved below. The precise extent to which their
combined formulation is new in the wider literature has not been established
by a comprehensive priority search. The finite verification program audits
examples and formulas; it is not a substitute for the proofs and is not a
Lean formalization.

\subsection{Four notions that must be distinguished}
An inverse may exist as a Hahn series without its successive approximations
converging in the full valuation topology. A finite list of coefficients
may be computable even when all coefficients below a chosen valuation
cutoff cannot be listed finitely. Finally, a formal Hahn series is not
automatically a germ of a real or complex function.

We therefore keep four questions separate: formal existence; strong Hahn
summation; valuation convergence; and analytic realization. The first three
are addressed under explicit hypotheses. Borel summability, resurgence,
Hardy-field realization, and a general natural derivation on a full
logarithmic--exponential transseries field are not consequences claimed here.

\section{Hahn series and the all-word support lemma}
\label{sec:words}

\subsection{Conventions}
Throughout, $\Gamma$ is a set-sized totally ordered abelian group, written
additively, and $K$ is a field of characteristic zero. The order is
translation invariant. We write $R=K[L]$, where $L$ is an indeterminate, and
\[
 \HH=R((t^\Gamma))
 =\left\{\sum_{\gamma\in T}p_\gamma(L)t^\gamma:
            T\subseteq\Gamma\text{ well ordered},\ p_\gamma\in R\right\}.
\]
Zero coefficients are omitted from the support. Our orientation is that
larger exponents represent smaller monomials. For $f\ne0$ let
$v(f)=\min\supp f$, and set $v(0)=+\infty$.

A family $(f_i)_{i\in I}$ is \emph{Hahn summable} if its union of supports is
well ordered and, for each exponent, only finitely many family members
have a nonzero coefficient there. Its sum is defined using those finite
coefficient sums. Convergence of an infinite sum of scalars does not make
an otherwise inadmissible family Hahn summable.

In this article, a sequence $f_n$ \emph{converges strongly in the Hahn sense}
to $f$ if its supports and that of $f$ lie in one well-ordered set and every
coefficient is eventually exactly the corresponding coefficient of $f$.
For such a sequence the successive-difference family is Hahn summable, and
\[
 f=f_0+\sum_{n\ge0}(f_{n+1}-f_n).
\]
This is the explicitly used sequential notion; it must not be replaced by
valuation convergence without proof.

\begin{lemma}[Binary support calculus]\label{lem:binary}
If $T,U\subseteq\Gamma$ are well ordered, then $T+U$ is well ordered and,
for every $\alpha\in\Gamma$, there are only finitely many pairs
$(\tau,u)\in T\times U$ satisfying $\tau+u=\alpha$.
\end{lemma}
\begin{proof}
Every sequence in a well-ordered set has an infinite nondecreasing
subsequence. Indeed, if some value occurs infinitely often use it; otherwise
successively choose a minimum of a remaining tail and pass beyond its last
occurrence. If a strictly decreasing sequence in $T+U$ existed, choose one
representation of each member and pass to a nondecreasing subsequence of
first coordinates. The second coordinates would strictly decrease, a
contradiction.

For infinitely many different representations of a fixed $\alpha$, the
first coordinates are different. A nondecreasing subsequence is then
strictly increasing, forcing the corresponding second coordinates to
strictly decrease. This proves finiteness of the fiber.
\end{proof}

It follows that the Cauchy product is well defined on $\HH$. It is an
integral domain, with $v(fg)=v(f)+v(g)$ for nonzero $f,g$. The following
stronger statement concerns words of \emph{unbounded length}, not merely a
fixed finite product.

\begin{lemma}[Positive all-word lemma]\label{lem:words}
Let $S\subseteq\Gamma_{>0}$ be well ordered. Its additive monoid
\[
 M(S)=\{0\}\cup\bigcup_{n\ge1}nS
\]
is well ordered. Each $\alpha\in\Gamma$ is the sum of only finitely many
ordered finite words over $S$. The empty word is permitted and has sum zero.
\end{lemma}
\begin{proof}
Order finite words by subsequence embedding with coordinatewise comparison:
$(a_1,\ldots,a_m)$ embeds in $(b_1,\ldots,b_n)$ if there are
$1\le j_1<\cdots<j_m\le n$ with $a_i\le b_{j_i}$.
The word theorem says that every infinite sequence of words has an earlier
word embedding in a later one. For completeness, its well-ordered-alphabet
case is proved here by the minimal-bad-sequence argument.

Suppose a bad sequence exists. Choose its first word of minimal possible
length among words beginning a bad sequence, then its second word of
minimal length among continuations of the chosen prefix, and continue.
No word is empty. Write $w_i=v_i a_i$, separating the last letter, and
choose indices $i_0<i_1<\cdots$ on which $a_{i_j}$ is nondecreasing.
The sequence
\[
 w_0,\ldots,w_{i_0-1},v_{i_0},v_{i_1},v_{i_2},\ldots
\]
cannot be bad, because its first new word is shorter than the word chosen
minimally at that position. An embedding among its old words contradicts
badness of the original sequence. An old word embedding in $v_{i_j}$ also
embeds in $w_{i_j}$. Finally, an embedding $v_{i_j}\preceq v_{i_k}$ extends,
using $a_{i_j}\le a_{i_k}$, to $w_{i_j}\preceq w_{i_k}$. Each possibility
is impossible. This proves the word theorem in the needed case.

Now an embedding of words makes their sums nondecreasing. Moreover, their
sums are \emph{strictly} increasing unless the two words are identical:
an omitted letter is positive, and a strictly larger matched letter also
strictly increases the sum. Therefore words representing a strictly
decreasing sequence of sums would be a bad sequence. This proves that
$M(S)$ is well ordered. Infinitely many different words with one fixed sum
would also be a bad sequence. This proves the finite-word assertion.
\end{proof}

The combinatorial mechanism is the classical Higman--Neumann support
principle \cite{Higman}. The proof makes clear why no real-valued lower bound
on word length is required.

% ed. (2026-09-29): the repository already has this lemma in Lean; the
% article's formalization plan (sec:formalization) lists it as a first layer.
\begin{ednote}\label{ed:words-lean}
Both assertions of Lemma~\ref{lem:words} are already formalized in the
repository's Lean library for surreal Hahn series, over an arbitrary
linearly ordered cancellative abelian group with no Archimedean
hypothesis \cite{EdSurrealLean}. The finite-word assertion is
\path{Surreal.HahnSeries.finite_words_of_sum_eq}
(\path{Algebra/SurrealNumbers/Surreal/HahnSeries/Neumann.lean}: all
finite lists over a positive partially well-ordered $S$ with a given sum
form a finite set, proved by Higman's lemma) and, in monoid form,
\path{Surreal.HahnSeries.finite_wordsWithSum}
(\path{Algebra/SurrealNumbers/Surreal/HahnSeries/NeumannWords.lean}).
Well-ordering of $M(S)$ is
\path{Surreal.HahnSeries.isPWO_closure_of_pos} in the same file, and
\path{Surreal.HahnSeries.neumann_positive} in \texttt{Neumann.lean}
combines it with finiteness of nondecreasing words. These declarations
were not visible to this article's targeted inspection, which read only
\texttt{Analysis/Transseries}; the ``All-word support'' layer of
Section~\ref{sec:formalization} can import them.
\end{ednote}

\begin{corollary}[Anchored words]\label{cor:anchored}
For well-ordered $T\subseteq\Gamma$ and positive well-ordered $S$, the set
$T+M(S)$ is well ordered. For each $\alpha$, there are only finitely many
pairs consisting of an anchor $\tau\in T$ and an ordered word over $S$
whose sum with $\tau$ is $\alpha$.
\end{corollary}
\begin{proof}
Apply Lemma~\ref{lem:binary} to $T$ and $M(S)$. There are finitely many
possible pairs of anchor and word sum, and Lemma~\ref{lem:words} supplies
finitely many words for each such sum.
\end{proof}

\begin{remark}[Where the Archimedean proof breaks]\label{rem:archbreak}
If $\Gamma\subseteq\R$, $\delta=\min S>0$, and a word has sum $\alpha$, then
its length is at most $\alpha/\delta$. In an arbitrary group, the inequalities
$n\delta\le\alpha$ may hold for every $n$. This invalidates that numerical
bound, not Lemma~\ref{lem:words}. Replacing the lemma by the failed bound
would discard valid inverses.
\end{remark}

\section{Support depth and a complete filtration}
\label{sec:filtration}

Assume from now on that $S$ is nonempty, positive, and well ordered, and put
$M=M(S)$. Define
\begin{equation}\label{eq:depth}
 \depth_S(\alpha)=\max\{n\in\N:\alpha\in nS\},\qquad \alpha\in M,
 \qquad 0S=\{0\}.
\end{equation}
Thus $\depth_S(0)=0$. This is the \emph{maximum} representation length, not
the minimum number of generators. Maximum length is the relevant invariant
when a coefficient may receive contributions from many nonlinear orders.
Lemma~\ref{lem:words} makes it finite.

For $r\ge0$ define
\begin{equation}\label{eq:Ir}
 S_{\ge r}=\bigcup_{n\ge r}nS,\qquad
 \AS=\{f\in\HH:\supp f\subseteq M\},\qquad
 I_r=\{f\in\AS:\supp f\subseteq S_{\ge r}\}.
\end{equation}
Here $I_0=\AS$, and $I_1$ is its augmentation ideal: it consists precisely
of the series with no exponent-zero block. We call $I_r$ the support-depth
filtration. We do not identify $I_r$ with an ordinary algebraic ideal power
without additional justification.

\begin{proposition}[Complete support filtration]\label{prop:complete}
The sets $I_r$ are ideals of $\AS$, and
\[
 I_{r+1}\subseteq I_r,\qquad I_pI_q\subseteq I_{p+q},\qquad
 \bigcap_{r\ge0}I_r=\{0\}.
\]
The ring $\AS$ is complete for this filtration. If $f_n\in I_{r_n}$ and
$r_n\to+\infty$ as integers, then $(f_n)$ is Hahn summable. Every element
of $1+I_1$ is a unit, with inverse given by its geometric series.
\end{proposition}
\begin{proof}
Concatenating words gives $S_{\ge p}+S_{\ge q}\subseteq S_{\ge p+q}$;
concatenating an arbitrary additional word shows that each $I_r$ is an
ideal. If $f$ belongs to every $I_r$, its coefficient at $\alpha$ must vanish
as soon as $r>\depth_S(\alpha)$, proving separatedness.

For the family $(f_n)$, the union of supports lies in the well-ordered set
$M$. A fixed $\alpha$ can occur only for indices with
$r_n\le\depth_S(\alpha)$, of which there are finitely many. Hence the family
is Hahn summable.

For completeness, let $(g_n)$ be filtration-Cauchy. For every $\alpha$, the
coefficients of $g_n$ at $\alpha$ eventually stabilize: use any filtration
index greater than $\depth_S(\alpha)$. Let $g$ consist of these stabilized
coefficients. Its support is a subset of $M$, so it lies in $\AS$. Fix $r$.
Beyond a Cauchy index, $g_n-g_m\in I_r$ for all $n,m$. Passing coefficientwise
to the stabilized values gives $g_n-g\in I_r$. This proves completeness.

Finally, if $a\in I_1$, then $a^n\in I_n$, so $\sum_{n\ge0}(-a)^n$ is
summable. The usual finite geometric identity, followed coefficientwise,
shows that this is $(1+a)^{-1}$.
\end{proof}

The same conclusions hold entrywise for finite vectors and matrices. In
particular, $\Id+J$ is invertible whenever every entry of $J$ is in $I_1$.

\begin{lemma}[Coefficient stopping rule]\label{lem:stop}
If $f-g\in I_r$ and $r>\depth_S(\alpha)$, then
$\coef{t^\alpha}f=\coef{t^\alpha}g$. For a finite nonempty target set
$E\subseteq M$, it suffices that
\[
 r>H_E,\qquad H_E=\max_{\alpha\in E}\depth_S(\alpha).
\]
\end{lemma}
\begin{proof}
No exponent of depth less than $r$ is in $S_{\ge r}$.
\end{proof}

Notice that this is a coefficient certificate, not a claim that every
initial segment of $M$ is finite. The distinction persists even when $S$
itself is finite.

\section{Small-coefficient implicit systems}
\label{sec:implicit}

The support mechanism is not restricted to one inverse equation. A finite
system makes the source of the extra Newton gain particularly transparent.

Fix $d\ge1$ and formal series in $d$ unknowns
\begin{equation}\label{eq:A}
 A_i(U)=\sum_{\nu\in\N^d}a_{i,\nu}U^\nu,
 \qquad \supp a_{i,\nu}\subseteq S,
 \qquad 1\le i\le d.
\end{equation}
There need not be a uniform bound on the degrees of the coefficient
polynomials in $L$. The coefficients need not be summable as $\nu$ varies
before evaluation. Their common positive support bound is the hypothesis.
Set
\[
 \Psi(U)=U+A(U).
\]

\begin{lemma}[Evaluation and nonlinear estimates]\label{lem:nonlinear}
For $U\in I_1^d$, the expressions $A(U)$ and all their formal partial
derivatives are defined by Hahn-summable families. Their values and all
entries of $A'(U)$ lie in $I_1$. For $U,V\in I_1^d$ with
$U-V\in I_r^d$, $r\ge1$, one has
\begin{align}
 A(U)-A(V)&\in I_{r+1}^d,\label{eq:lipschitz}\\
 A(U)-A(V)-A'(V)(U-V)&\in I_{2r+1}^d.\label{eq:quadratic}
\end{align}
\end{lemma}
\begin{proof}
A term with $|\nu|=m$ in $A(U)$ belongs to $I_{m+1}$. There are finitely
many multi-indices of each total degree because $d$ is finite. The common
support is contained in $M$, and a fixed exponent permits only finitely many
such degrees by its finite depth. Each fixed-degree expression is an
ordinary finite sum of Hahn products. This proves summability. The same
argument applies after any fixed number of partial differentiations;
the coefficient still supplies at least one support letter.

For each monomial, $U^\nu-V^\nu$ is a sum of terms with at least one factor
from $U-V$, and the associated coefficient supplies a letter from $S$.
This proves \eqref{eq:lipschitz}. Subtracting the linear part leaves at
least two factors from $U-V$, so its support depth is at least $2r+1$.
All expansions are legitimate by the preceding fixed-coefficient finiteness
argument.
\end{proof}

\begin{theorem}[Implicit solution and sharp Newton estimate]
\label{thm:implicit}
Under \eqref{eq:A}, the equation $\Psi(U)=0$ has a unique solution
$U_*\in I_1^d$. Picard iteration
\[
 P_0=0,\qquad P_{n+1}=-A(P_n)
\]
converges strongly in the Hahn sense to $U_*$, with
\begin{equation}\label{eq:picarderr}
 P_n-U_*\in I_{n+1}^d.
\end{equation}
Newton iteration
\begin{equation}\label{eq:newton}
 N_0=0,\qquad
 N_{n+1}=N_n-[\Id+A'(N_n)]^{-1}\Psi(N_n)
\end{equation}
is well defined and converges strongly to the same solution. Its error
satisfies
\begin{equation}\label{eq:newtondepth}
 N_n-U_*\in I_{r_n}^d,\qquad r_n=2^{n+1}-1.
\end{equation}
More generally, one Newton step sends an error in $I_r^d$ into
$I_{2r+1}^d$.
\end{theorem}
\begin{proof}
Picard iteration stays in $I_1^d$. Its first increment belongs to $I_1^d$,
and \eqref{eq:lipschitz} shows inductively that
$P_{n+1}-P_n\in I_{n+1}^d$. Proposition~\ref{prop:complete} gives a strong
limit $U_*\in I_1^d$. The same Lipschitz estimate permits passage to the
limit in $P_{n+1}=-A(P_n)$, yielding $\Psi(U_*)=0$.
The tail of the increment family gives \eqref{eq:picarderr}.

If $U,V\in I_1^d$ are solutions, their difference first belongs to $I_1^d$.
Equation \eqref{eq:lipschitz} and the fixed-point equation then put it in
$I_2^d$, then in $I_3^d$, and so forth. Separatedness makes it zero.
No valuation cofinality enters this argument.

For Newton iteration, $J(U)=\Id+A'(U)$ is a matrix over $\AS$ congruent to
$\Id$ modulo $I_1$, hence is invertible by its matrix geometric series.
The correction belongs to $I_1^d$, so the iteration stays in that domain.
Suppose $e=U-U_*\in I_r^d$. Taylor expansion about $U$, evaluated at $U_*$,
gives
\[
 0=\Psi(U_*)=\Psi(U)-J(U)e+R,
 \qquad R\in I_{2r+1}^d.
\]
The linear part $U$ of $\Psi$ has no quadratic remainder; all higher
terms come from $A$, whose coefficients carry the additional support
letter. Consequently,
\[
 U-J(U)^{-1}\Psi(U)-U_*=J(U)^{-1}R\in I_{2r+1}^d.
\]
Initially $N_0-U_*\in I_1^d$. The recurrence
$r_{n+1}=2r_n+1$, $r_0=1$, has solution $r_n=2^{n+1}-1$.
The resulting errors and increments form Hahn-summable families, and each
coefficient eventually agrees with $U_*$. This proves strong convergence.
\end{proof}

% ed. (2026-09-29): credits the identical rate of the predecessor article.
\begin{ednote}\label{ed:newton-rate}
Applied to reversion (Section~\ref{sec:sharp}), where the normalized
correction is $U=t^hC$, with $\Gamma\subseteq\R$ and $h=1$,
\eqref{eq:newtondepth} gives $v(C_n-C_*)\ge(2^{n+1}-1)\delta-1$,
$\delta=\min S$, and the step
$r\mapsto2r+1$ becomes $e\mapsto2e+\alpha+2$ with $\delta=\alpha+1$: this is
equation \texttt{eq:newtonrate} of Theorem~4.1 (\texttt{thm:newton}) in
\path{Analysis/Transseries/docs/series-and-transseries/Support_Controlled_Reversion_One_Exponential/reversion_and_one_exponential.tex},
which this article does not cite. The depth formulation extends that rate
to groups in which the multiples of $\delta$ are not cofinal;
Theorem~\ref{thm:sharp} shows it cannot be improved.
\end{ednote}

\begin{corollary}[Exact Newton round count]\label{cor:rounds}
For a finite nonempty set of requested exponents $E\subseteq M$, all those
coefficients of $N_n$ are exact whenever
\begin{equation}\label{eq:rounds}
 2^{n+1}-1>H_E.
\end{equation}
In particular,
\[
 n=\left\lfloor\log_2(H_E+1)\right\rfloor
\]
is sufficient, including $n=0$ when $H_E=0$.
\end{corollary}
\begin{proof}
Combine \eqref{eq:newtondepth} with Lemma~\ref{lem:stop}. The displayed
integer is the least nonnegative $n$ satisfying \eqref{eq:rounds}.
\end{proof}

\begin{proposition}[Uniqueness beyond the generated monoid]
\label{prop:allpositive}
The solution in Theorem~\ref{thm:implicit} is the only solution whose
components are arbitrary positive-support Hahn series, even if their
supports were not initially assumed to lie in $M(S)$.
\end{proposition}
\begin{proof}
For any finite collection of positive-support inputs, their combined
support is well ordered and positive. The anchored-word lemma justifies
evaluation of \eqref{eq:A} on these inputs. Write $\delta=\min S$.
If $U,V$ are such inputs, and
$\epsilon=\min_i v(U_i-V_i)$ is finite, the monomial difference
factorization gives
\[
 v\bigl(A_i(U)-A_i(V)\bigr)\ge\epsilon+\delta
 \quad\text{for every }i.
\]
All other factors have nonnegative valuations, and every coefficient of
$A$ has valuation at least $\delta$. If both inputs were solutions, some
component of $U-V$ would have valuation $\epsilon$, whereas its equal
negative feedback difference would have strictly larger valuation.
This is impossible. A \emph{single} strict gain suffices for uniqueness;
its multiples need not be cofinal.
\end{proof}

\section{An explicit finite-tree coefficient formula}
\label{sec:trees}

The fixed-point proof has a useful nonrecursive interpretation. Write
\[
 a_{i,\nu}=\sum_{s\in S}a_{i,\nu,s}(L)t^s.
\]
Consider finite rooted trees whose vertices have types in
$\{1,\ldots,d\}$. A vertex of type $i$ has $\nu_j$ ordered child slots of
type $j$, grouped in increasing order of $j$, for some $\nu\in\N^d$.
Give this vertex an additional label $s\in S$ and weight
$-a_{i,\nu,s}(L)$. Zero-weight decorations may simply be omitted.
The exponent of a decorated tree is the sum of its vertex labels; its
coefficient weight is the product of its vertex weights.

\begin{theorem}[Finite tree formula]\label{thm:trees}
For $\alpha\in M$ and root type $i$,
\begin{equation}\label{eq:treeformula}
 \coef{t^\alpha}U_{*,i}
 =\sum_{\substack{\mathcal T\text{ decorated tree of root type }i\\
                    \sum_{v\in\mathcal T}s_v=\alpha}}
       \prod_{v\in\mathcal T}\bigl(-a_{i_v,\nu_v,s_v}(L)\bigr).
\end{equation}
This is a finite sum, and every contributing tree has at most
$\depth_S(\alpha)$ vertices. There is no automorphism-group denominator:
child slots are ordered as specified above.
\end{theorem}
\begin{proof}
List the vertices in depth-first order. Their labels form an ordered word
over $S$ with sum $\alpha$. There are finitely many such words and a finite
upper bound on their lengths. For a fixed number of vertices there are
finitely many rooted ordered tree shapes, finitely many vertex types, and
finitely many possible child-count vectors. Thus the displayed family is
finite at each exponent and has support in $M$.

Decompose a tree at its root. The root supplies one coefficient of $-A_i$;
the ordered subtrees in its child slots supply the factors
$U_1^{\nu_1}\cdots U_d^{\nu_d}$. This is exactly the fixed-point equation.
Conversely, expanding a product in that equation independently chooses a
subtree for each slot, producing every decorated tree with precisely its
multiplicity. Hence the Hahn sum over all trees is a solution in $I_1^d$.
Uniqueness proves \eqref{eq:treeformula}.
\end{proof}

The tree formula is not presented as a new general theory of implicit
functions. Its role here is to exhibit a finite certificate for an
individual coefficient: the dependence on arbitrarily high nonlinear
orders disappears once the target exponent is fixed. Picard iteration
organizes the same trees by height. The support-depth bound instead
controls their number of vertices, and remains useful when infinitely many
support exponents lie below a valuation cutoff.

\section{The fixed-shift power--logarithmic differential algebra}
\label{sec:differential}

\subsection{The necessary additional data}
In a real exponent group, the exponent $\gamma$ can also be used as a scalar
coefficient. This is not meaningful in an arbitrary ordered abelian group
without a specified map. We therefore fix
\begin{equation}\label{eq:character}
 h\in\Gamma_{>0},\qquad \chi:\Gamma\longrightarrow(K,+)
 \text{ additive},\qquad \chi(h)=1.
\end{equation}
No order compatibility of $\chi$ is assumed. Put $X=t^{-h}$, and define
\begin{equation}\label{eq:D}
 D\bigl(t^\gamma p(L)\bigr)
 =t^{\gamma+h}\bigl(p'(L)-\chi(\gamma)p(L)\bigr).
\end{equation}
Then $DX=1$ and $DL=t^h$. When $\Gamma\subseteq\R$, $K=\R$, $h=1$ and
$\chi(\gamma)=\gamma$, this is exactly the power--logarithmic derivative
used by the repository source.

The data \eqref{eq:character} can always be supplied after a suitable choice
of scalar map: an ordered abelian group is torsion free, hence embeds in
$\Gamma\otimes_\Z\Q$; extend $h$ to a rational vector-space basis and send
it to $1$. Sending the other basis elements to zero gives a map into $\Q$.
If an injective additive map is desired, send them instead to algebraically
independent indeterminates and take a corresponding rational-function
field. These choices need not be canonical and need not provide an
analytic interpretation.

\begin{proposition}[Derivative and Euler operator]\label{prop:derivative}
Equation \eqref{eq:D} extends to a strongly summation-preserving derivation
of $\HH$. For every $j\ge0$,
\begin{equation}\label{eq:Diter}
 D^j\bigl(t^\gamma p(L)\bigr)
 =t^{\gamma+jh}\prod_{i=0}^{j-1}
       (\partial_L-\chi(\gamma)-i)p(L).
\end{equation}
In particular $v(D^j f)\ge v(f)+jh$. The Euler operator
\begin{equation}\label{eq:Euler}
 \nabla=XD=t^{-h}D
\end{equation}
is a derivation preserving each support exponent, with
$\nabla(t^\gamma p)=t^\gamma(p'-\chi(\gamma)p)$.
\end{proposition}
\begin{proof}
Translation by $h$ preserves well ordering and finite support fibers, so
termwise application is strong. The product rule follows on two monomial
blocks from additivity of $\chi$ and the ordinary polynomial product rule;
strong summation extends it to all series. Induction on $j$, using
$\chi(\gamma+jh)=\chi(\gamma)+j$, gives \eqref{eq:Diter}.
Multiplication of a derivation by an element of a commutative ring is again
a derivation; multiplying \eqref{eq:D} by $t^{-h}$ gives the last formula.
\end{proof}

\begin{remark}[The scope is deliberate]
The ordered exponent group controls support; the character $\chi$ controls
the differential weights. A character with kernel gives additional
constant monomials. Nothing in the reversion proof requires the constant
field to equal $K$. A full transseries field can have a more complicated
derivative whose support shifts vary with the monomial. Such a derivative
is outside the fixed-shift theorem and is a further research direction,
not silently covered by \eqref{eq:D}.
\end{remark}

\section{Strong Taylor substitution without cofinality}
\label{sec:substitution}

Let $B\in\HH$ satisfy $v(B)>-h$, and put
\[
 b=t^hB,\qquad H_B=X+B.
\]
The support $S=\supp b$ is positive and well ordered. The zero case $B=0$
is understood separately as the identity substitution. Define
\begin{equation}\label{eq:Taylor}
 T_Bf=f\circ H_B:=\sum_{j\ge0}\frac{B^j}{j!}D^jf.
\end{equation}
Here the notation $\circ$ denotes the substitution being constructed,
not an assumed analytic composition operation.

\begin{theorem}[Strong substitution calculus]\label{thm:substitution}
For every $f\in\HH$, formula \eqref{eq:Taylor} is Hahn summable. If
$T=\supp f$, its support is contained in $T+M(S)$. The map $T_B$ is a
strongly summation-preserving $K$-algebra homomorphism. It acts by
\begin{equation}\label{eq:generators}
 T_B(t^\gamma)=t^\gamma(1+b)^{-\chi(\gamma)},\qquad
 T_B(L)=L+\log(1+b).
\end{equation}
It preserves the leading coefficient block and valuation of every nonzero
series. Moreover,
\begin{align}
 D(T_Bf)&=(1+DB)T_B(Df),\label{eq:chain}\\
 T_C(T_Bf)&=T_{C+T_CB}f\label{eq:association}
\end{align}
whenever $v(C)>-h$.
\end{theorem}
\begin{proof}
By \eqref{eq:Diter}, the $j$th Taylor summand has support in
\[
 j(\supp B)+(T+jh)=T+jS.
\]
The anchored-word lemma proves that the union is well ordered and that a
fixed target receives contributions from only finitely many Taylor orders
and finitely many input anchors. This proves summability. For an arbitrary
Hahn-summable input family, its combined anchors form a well-ordered set
and each anchor occurs in only finitely many members. The same argument
therefore permits both orders of summation and proves strong preservation.

For $t^\gamma$, the iterated derivative formula gives the binomial series
in \eqref{eq:generators}. For $L$, its positive derivatives are
\[
 D^jL=(-1)^{j-1}(j-1)!t^{jh}\quad(j\ge1),
\]
which give the logarithm. The binomial and logarithmic series are themselves
summable because $b$ has positive support. On monomial blocks,
\eqref{eq:generators} respects products, by additivity of $\chi$ and the
identity $(1+b)^{a+c}=(1+b)^a(1+b)^c$. It respects polynomial operations in
$L$ as well. These formal identities may be checked in one ordinary
power-series variable and specialized using the word lemma. Strong
summation then proves the homomorphism assertion on $\HH$.

Every positive Taylor order shifts an anchor by a positive support word.
It therefore contributes nothing to the least exponent of $f$.
The leading block is unchanged.

For the chain rule, differentiating the summable Taylor family gives
\[
 \begin{split}
 D(T_Bf)
 &=\sum_{j\ge0}\frac{B^j}{j!}D^{j+1}f
   +(DB)\sum_{j\ge1}\frac{B^{j-1}}{(j-1)!}D^jf\\
 &=(1+DB)T_B(Df).
 \end{split}
\]
The derivative is strong and products of the relevant families are
summable, justifying the displayed manipulation.

For associativity, let $c=t^hC$ and set $E=C+T_CB$.
Since $T_C$ preserves valuation, $v(E)>-h$. Also
\[
 t^hE=c+(1+c)T_Cb,
 \qquad 1+t^hE=(1+c)(1+T_Cb).
\]
Applying $T_C$ to \eqref{eq:generators}, and using this product identity,
shows that $T_CT_B$ and $T_E$ agree on every $t^\gamma$ and on $L$.
For $L$ use $\log(uv)=\log u+\log v$ for commuting elements of $1+I_1$
in the monoid generated by the combined positive supports. Both sides
preserve Hahn sums, so they agree on every series. This proves
\eqref{eq:association}.
\end{proof}

\begin{proposition}[A difference estimate requiring only one gain]
\label{prop:difference}
If $v(B)>-h$ and $v(C),v(E)>-h$, then
\begin{equation}\label{eq:diffv}
 v(T_CB-T_EB)\ge v(C-E)+v(B)+h.
\end{equation}
In particular, $H_B\circ H_C=H_B\circ H_E$ implies $C=E$.
\end{proposition}
\begin{proof}
Subtract Taylor expansions and use
$C^j-E^j=(C-E)\sum_{a=0}^{j-1}C^aE^{j-1-a}$ for $j\ge1$.
Each resulting term has valuation at least
\[
 v(C-E)+v(B)+h
 +(j-1)\bigl(h+\min\{v(C),v(E)\}\bigr).
\]
The last summand is nonnegative. All families are summable by applying the
word lemma to the union of the normalized input supports. This proves
\eqref{eq:diffv}. If $C-E=-(T_CB-T_EB)$ and $C\ne E$, the right side would
have strictly larger valuation than the left, because $v(B)+h>0$.
\end{proof}

\section{Reversion and the ordinary Lagrange index}
\label{sec:reversion}

\subsection{Reduction to a small-coefficient equation}
Write the normalized input as
\begin{equation}\label{eq:binput}
 b=t^hB=\sum_{s\in S}b_s(L)t^s.
\end{equation}
Seek $G=X+C$ and put $U=t^hC$. By \eqref{eq:generators}, the equation
$H_B\circ G=X$ is equivalent to
\begin{equation}\label{eq:normalinverse}
 \Psi_B(U)=U+A_B(U)=0,
\end{equation}
where
\begin{equation}\label{eq:AB}
 A_B(U)=\sum_{s\in S}t^s(1+U)^{1-\chi(s)}
                    b_s\bigl(L+\log(1+U)\bigr).
\end{equation}
Indeed, a block of $B$ has exponent $s-h$ and hence differential weight
$\chi(s)-1$. The formal coefficient of every power of $U$ in $A_B$ has
support contained in $S$. Thus \eqref{eq:AB} is a scalar instance of
\eqref{eq:A}.

\begin{theorem}[Non-Archimedean power--logarithmic reversion]
\label{thm:reversion}
Assume \eqref{eq:character}. Every $H_B=X+B$ with $v(B)>-h$ has a unique
two-sided inverse of the form $G=X+C$, $v(C)>-h$. Its normalized correction
$U=t^hC$ belongs to $I_1$ for $S=\supp(t^hB)$. In particular, its support
is contained in $M(S)\setminus\{0\}$.

The inverse is given by the ordinary positive-integer Lagrange sum
\begin{equation}\label{eq:LB}
 G=X+\sum_{n\ge1}\frac{(-1)^n}{n!}D^{n-1}(B^n).
\end{equation}
The normalized $n$th summand has support in $nS$. Consequently, the
coefficient at exponent $\alpha$ of $U$ uses at most the outer orders
$1\le n\le\depth_S(\alpha)$, whether or not $n\min S$ is cofinal in
$\Gamma$.
\end{theorem}
\begin{proof}
Theorem~\ref{thm:implicit}, applied to \eqref{eq:AB}, gives a solution
$U\in I_1$, hence a right substitution inverse $G=X+t^{-h}U$ with
$H_B\circ G=X$. It lies in the same near-identity class. Apply the same
existence argument to $G$ to obtain $H$ with $G\circ H=X$.
Associativity from Theorem~\ref{thm:substitution} gives
\[
 H_B=H_B\circ(G\circ H)=(H_B\circ G)\circ H=H.
\]
Therefore $G\circ H_B=X$, proving two-sided inversion.
Proposition~\ref{prop:difference} proves uniqueness among all near-identity
corrections, not only those with a prespecified generated support.

To derive \eqref{eq:LB}, introduce an independent ordinary formal variable
$z$ and solve
\[
 C(z)=-z\,\Phi(C(z)),\qquad
 \Phi(y)=\sum_{j\ge0}\frac{D^jB}{j!}y^j
 \quad\text{in }\HH[[y]].
\]
Classical coefficient Lagrange inversion, over any commutative
$\Q$-algebra, gives
\[
 [z^n]C(z)=\frac{(-1)^n}{n}[y^{n-1}]\Phi(y)^n.
\]
A proof of the exact form used here is included in
Appendix~\ref{app:lagrange}. The formal Taylor map
$f\mapsto\sum_jD^jf\,y^j/j!$ is a ring homomorphism, by the iterated
Leibniz formula. Hence
\[
 [y^{n-1}]\Phi(y)^n=\frac{1}{(n-1)!}D^{n-1}(B^n),
\]
which proves the coefficient formula with $z$ retained.

It remains essential to justify setting $z=1$. The normalized $n$th
coefficient has support in $nS$: $B^n$ shifts $nS$ by $-nh$,
$D^{n-1}$ shifts it by $(n-1)h$, and multiplication by $t^h$ removes the
remaining shift. Lemma~\ref{lem:words} therefore makes the normalized
coefficient family Hahn summable. The parameter equation can be written
\[
 U(z)+zA_B(U(z))=0,\qquad U(z)=t^hC(z).
\]
Upon fully expanding a coefficient of this equation, each occurrence of
$z$ is accompanied by one input support letter. For any target exponent
there are only finitely many decorated words, so substitution $z=1$
commutes with every sum and product in the equation. The resulting Hahn
sum solves \eqref{eq:normalinverse}; uniqueness identifies it with $U$.
Finally, no normalized $n$th term can contribute at $\alpha$ for
$n>\depth_S(\alpha)$.
\end{proof}

% ed. (2026-09-29): names the other proofs of the same Lagrange sum.
\begin{ednote}\label{ed:reversion-other-proofs}
At $\Gamma\subseteq\R$, $h=1$, $\chi=\Id$, \eqref{eq:LB} is Theorem~3.2
(\texttt{thm:LB}, equation \texttt{eq:LB}) of the predecessor reversion
article
(\path{Analysis/Transseries/docs/series-and-transseries/Support_Controlled_Reversion_One_Exponential/reversion_and_one_exponential.tex}),
whose $A$ and $B$ are the $B$ and $C$ here, and Corollary~\ref{cor:EulerLB}
agrees term by term with its block formula \texttt{eq:blockLB}. Over a
non-Archimedean exponent group, existence, uniqueness and the same sum
without logarithmic blocks are Lemma \texttt{e:lem-reversion} of the
surcomplex analysis volume
(\path{Algebra/SurrealNumbers/docs/surcomplex/analysis/article.tex}),
proved by the same auxiliary-parameter specialization. What this theorem
adds is the power--logarithmic block calculus over a non-Archimedean group
with fixed shift data $(h,\chi)$; its outer-order bound
$n\le\depth_S(\alpha)$ replaces the real-valued bound \texttt{eq:rbound}
of the predecessor.
\end{ednote}

\subsection{A normalized Euler formula}
The derivative formula can be kept entirely inside the generated support
algebra, an advantage for finite quotient computation.

\begin{corollary}[Euler-operator reversion]\label{cor:EulerLB}
With $b=t^hB$ and $\nabla=XD$, the normalized inverse is
\begin{equation}\label{eq:EulerLB}
 \boxed{\quad
 U=\sum_{n\ge1}\frac{(-1)^n}{n!}
       \left(\prod_{j=2}^{n}(\nabla+j)\right)b^n.
 \quad}
\end{equation}
The empty product at $n=1$ is the identity. Each factor preserves support,
and the $n$th term is supported in $nS$.
\end{corollary}
\begin{proof}
Since $B^n=t^{-nh}b^n$, repeated use of
$D(t^\gamma f)=t^{\gamma+h}(\nabla f-\chi(\gamma)f)$ gives
\[
 t^hD^{n-1}(B^n)
  =(\nabla+2)(\nabla+3)\cdots(\nabla+n)b^n.
\]
The factors commute because each is a polynomial in the same operator.
Insert this identity in \eqref{eq:LB}.
\end{proof}

For example, the first three normalized orders are
\[
 U=-b+\tfrac12(\nabla+2)b^2
      -\tfrac16(\nabla+2)(\nabla+3)b^3+\cdots.
\]
These are nonlinear support orders. With resonances, several such orders
can contribute to the same exponent; the formula neither chooses a free
multi-index representation nor ignores these identifications.

\subsection{Support and exponent-lattice stability}

\begin{corollary}[No new exponent lattice]\label{cor:lattice}
If $S$ lies in a subgroup $\Lambda\subseteq\Gamma$, then the normalized
inverse support lies in $\Lambda$. Every normalized Picard and Newton
iterate also has support in $M(S)$. If $h\in\Lambda$, the unnormalized
corrections remain in $\Lambda$ as well.
\end{corollary}
\begin{proof}
The monoid $M(S)$ is contained in $\Lambda$, and every construction takes
place in $\AS$. Multiplication by $t^{-h}$ shifts exponents by $-h$.
\end{proof}

Thus the denominator-stability conclusion in a fixed Puiseux lattice
extends to arbitrary ordered-group sublattices with the stated shift
qualification. This is a support assertion. It supplies neither an
algorithm for an unpresented group nor an analytic realization of its
monomials.

\section{Sharp Newton depth and the exact valuation boundary}
\label{sec:sharp}

\subsection{The Newton iteration attached to reversion}
Differentiating \eqref{eq:AB} with respect to $U$ gives
\begin{equation}\label{eq:ABderiv}
 A'_B(U)=\sum_{s\in S}t^s(1+U)^{-\chi(s)}
 \left[b'_s\bigl(L+\log(1+U)\bigr)
 +(1-\chi(s))b_s\bigl(L+\log(1+U)\bigr)\right].
\end{equation}
Equivalently, if $C=t^{-h}U$, then $A'_B(U)=T_C(DB)$.
Thus the Newton iteration of Theorem~\ref{thm:implicit} is exactly
\begin{equation}\label{eq:NewtonC}
 C_0=0,\qquad
 C_{n+1}=C_n-\frac{C_n+T_{C_n}B}{1+T_{C_n}(DB)}.
\end{equation}
The denominator is $1+$ a positive-support series, so it is a unit.
The normalized error $t^h(C_n-C)$ lies in $I_{2^{n+1}-1}$.

\subsection{A scalar model attaining every support bound}

\begin{theorem}[Sharpness at every Newton stage]\label{thm:sharp}
Let $z$ be an ordinary formal variable and consider
\begin{equation}\label{eq:scalar}
 \Psi(u)=u+\frac{z}{1+u}=0.
\end{equation}
Its zero-constant solution is
\begin{equation}\label{eq:Catalan}
 u_* =\frac{-1+\sqrt{1-4z}}2
      =-\sum_{m\ge1}\frac1m\binom{2m-2}{m-1}z^m.
\end{equation}
Newton iteration for the rational function \eqref{eq:scalar}, starting at
$u_0=0$, satisfies
\begin{equation}\label{eq:sharpseries}
 u_n-u_*=z^{r_n}+\Ocal(z^{r_n+1}),\qquad r_n=2^{n+1}-1.
\end{equation}
Consequently, the uniform support-depth bound in
Theorem~\ref{thm:implicit} cannot be increased by even one at any stage.
\end{theorem}
\begin{proof}
The root formula follows from $u_*^2+u_*+z=0$ and the branch with constant
term zero. Its coefficient expansion is obtained by the formal binomial
series.

Write $e=u-u_*$. A direct algebraic simplification of the Newton update,
using $z=-u_*(1+u_*)$, gives the exact error identity
\begin{equation}\label{eq:exactNewtonError}
 e_{\rm new}
 =\frac{-u_*e^2}{(1+u)^2-z}.
\end{equation}
The denominator has constant term $1$, and $-u_*=z+\Ocal(z^2)$.
Thus an error $e=z^r+\Ocal(z^{r+1})$ becomes
$z^{2r+1}+\Ocal(z^{2r+2})$, with leading coefficient still $1$.
Since $e_0=-u_*=z+\Ocal(z^2)$, induction proves
\eqref{eq:sharpseries}. Take the support set to consist of the single
exponent of $z$; its depth is ordinary power degree. Every asserted bound
is then attained.
\end{proof}

\begin{remark}[Do not change the equation before comparing Newton rates]
The rational equation $u+z/(1+u)=0$ and the polynomial equation
$u^2+u+z=0$ have the same small root but different Newton maps.
Multiplying an equation by a nonconstant unit need not preserve Newton
iterates. The support gain $2r+1$ and the verification program concern
\eqref{eq:scalar}, which is the equation arising from normalized reversion.
\end{remark}

\subsection{A complete criterion for uniform valuation convergence}
The valuation topology on $\HH$ has zero-neighborhoods
$\{f:v(f)>\beta\}$ for $\beta\in\Gamma$. For finite vectors use the
minimum component valuation. Say that $\delta>0$ is \emph{cofinal by
multiples} if, for every $\beta\in\Gamma$, some $n\in\N$ satisfies
$n\delta>\beta$.

\begin{theorem}[Exact uniform valuation criterion]
\label{thm:valuationcriterion}
Fix a nonempty positive well-ordered support set $S$, and let
$\delta=\min S$. The following are equivalent.
\begin{enumerate}[label=\textup{(\roman*)}]
\item The multiples of $\delta$ are cofinal in $\Gamma$.
\item For every finite-dimensional small-coefficient system
\eqref{eq:A} with coefficient supports contained in $S$, its Newton
sequence starting at zero converges in the full valuation topology.
\item For every such system, its Picard sequence starting at zero
converges in the full valuation topology.
\end{enumerate}
Strong Hahn convergence holds for every one of these systems even when
these equivalent conditions fail.
\end{theorem}
\begin{proof}
Every exponent in $S_{\ge r}$ is at least $r\delta$. Therefore
\eqref{eq:picarderr} and \eqref{eq:newtondepth} imply valuation convergence
when $n\delta$ is cofinal. This proves (i)$\Rightarrow$(ii),(iii).

Conversely, put $z=t^\delta$ and use the scalar system
$u+z/(1+u)=0$. Its feedback power-series coefficients are $(-1)^jz$,
all supported on $\{\delta\}\subseteq S$, so it is in the stated class.
By Theorem~\ref{thm:sharp}, its Newton error has valuation $r_n\delta$.
If the multiples of $\delta$ are not cofinal, choose $\beta$ bounding them
above. None of these errors has valuation greater than $\beta$; hence the
Newton sequence does not converge in the valuation topology.

For Picard iteration $p_{n+1}=-z/(1+p_n)$, subtraction of the fixed-point
equation gives
\[
 p_{n+1}-u_*
 =\frac{z(p_n-u_*)}{(1+p_n)(1+u_*)}.
\]
Starting from $p_0-u_*=z+\Ocal(z^2)$, this implies
$p_n-u_*=z^{n+1}+\Ocal(z^{n+2})$. The same bound $\beta$ prevents valuation
convergence. Thus (ii) and (iii) each imply (i). Strong convergence follows
from Theorem~\ref{thm:implicit} without any cofinality assumption.
\end{proof}

The quantifier ``for every system'' is essential. A particular equation
can have a finite exact iteration even if $\delta$ is noncofinal; for
example, $u+a=0$ is solved by one Newton update. The theorem classifies the
uniform guarantee supplied by the support class, not the behavior of every
individual equation separately.

\subsection{An actual reversion that is not valuation-Cauchy}
Take $\Gamma=\Z^2$ in lexicographic order and set
\[
 h=(0,1),\qquad \eta=(1,0).
\]
Then $nh<\eta$ for every integer $n$. One may take
$K=\Q(\omega)$ with $\omega$ an indeterminate and
$\chi(a,b)=a\omega+b$, so the character can be injective; the following
example does not depend on a large kernel of $\chi$.

Let
\begin{equation}\label{eq:ranktwoF}
 F=X+X^{-1},\qquad X=t^{-h},\qquad z=t^{2h}.
\end{equation}
Its normalized inverse equation is exactly \eqref{eq:scalar}, with
$G=X(1+u_*)$. Hence
\begin{equation}\label{eq:ranktwoinverse}
 G=\frac{X+\sqrt{X^2-4}}2,
\end{equation}
where the square root means $X\sqrt{1-4X^{-2}}$ in the near-identity branch.
The Newton corrections satisfy
\begin{equation}\label{eq:ranktwov}
 v(C_n-C_*)=(2r_n-1)h=(2^{n+2}-3)h<\eta.
\end{equation}
They are not even valuation-Cauchy: the leading exponent of
$C_{n+1}-C_n$ is the leading exponent of $C_n-C_*$, since the error at the
next stage has strictly larger exponent. Nevertheless, they converge
strongly to the exact inverse, and any specified coefficient becomes exact
after finitely many stages.

The formal closed expression \eqref{eq:ranktwoinverse} is perfectly finite.
What fails is a particular topological convergence requirement and the
possibility of listing the infinitely many nonzero terms below the cutoff
$\eta$ in a finite monomial list. These are not failures of the inverse.

\section{When is an entire valuation cutoff finite?}
\label{sec:cuts}

The support criterion for a finite coefficient request is automatic, but
that for a whole valuation cutoff is not.

\begin{theorem}[Exact cut-finiteness criterion]\label{thm:cut}
Let $S\subseteq\Gamma_{>0}$ be nonempty and well ordered, let
$\delta=\min S$, and let $\beta\ge0$. Then
\begin{equation}\label{eq:cutiff}
 M(S)\cap(-\infty,\beta]\text{ is finite}
\end{equation}
if and only if both conditions hold:
\[
 S\cap(-\infty,\beta]\text{ is finite},
 \qquad N\delta>\beta\text{ for some }N\ge1.
\]
\end{theorem}
\begin{proof}
If \eqref{eq:cutiff} holds, its subset of generators is finite.
Also, the distinct elements $n\delta$ belong to $M(S)$. If every
$n\delta\le\beta$, the cutoff would contain infinitely many of them,
a contradiction.

Conversely, suppose the two conditions hold. A word of sum at most
$\beta$ has length less than $N$, because each letter is at least
$\delta$. Every letter of such a word is at most $\beta$, because all
letters are positive. There are only finitely many permitted letters and
finitely many permitted lengths, hence only finitely many words and sums.
\end{proof}

For $\beta<0$ the intersection is empty. For finite $S$, only the
cofinality-across-that-cut condition remains. In the rank-two example,
$S=\{2h\}$ is finite but no multiple of $2h$ crosses $\eta$, so the cutoff
is infinite.

\begin{remark}[Existential finiteness is not an input representation]
For an arbitrary infinite support supplied with no effective description,
Lemma~\ref{lem:words} proves that a coefficient has finitely many
contributors; it does not provide an algorithm to discover them. Effective
computation requires a presentation of the support and exponent arithmetic,
or a finite certificate that its relevant part is complete. The next
section supplies such certificates for finite generator lists, and
Section~\ref{sec:infinite} gives an explicit infinite-support family.
\end{remark}

\section{Minimal finite coefficient quotients and certificates}
\label{sec:certificates}

\subsection{The exact dependence set}
Let $E\subseteq M(S)$ be finite and nonempty. Define its monoid-divisor
closure by
\begin{equation}\label{eq:divisors}
 \Div_S(E)=\{\beta\in M(S):\alpha-\beta\in M(S)
                              \text{ for some }\alpha\in E\}.
\end{equation}
The notation involves group subtraction, but both summands in a permitted
decomposition must lie in the positive monoid. This is not an order
interval: a smaller exponent need not be a monoid divisor.

\begin{lemma}[Finite dependence closure]\label{lem:divisors}
The set $\Div_S(E)$ is finite, contains $0$ and $E$, and is downward closed
under monoid divisibility. Every $\beta\in\Div_S(E)$ satisfies
$\depth_S(\beta)\le H_E$.
\end{lemma}
\begin{proof}
For each target $\alpha$, Lemma~\ref{lem:binary} gives finitely many pairs
$(\beta,\gamma)\in M\times M$ with $\beta+\gamma=\alpha$.
Taking the finite union over $E$ proves finiteness. The containment
assertions follow by using zero as a summand. If $\beta$ divides $\alpha$
and $\beta'$ divides $\beta$, then $\beta'$ divides $\alpha$, proving
downward closure. A maximum-length word for $\beta$ can be concatenated
with a word for $\alpha-\beta$; its length is at most
$\depth_S(\alpha)\le H_E$.
\end{proof}

Put $D_E=\Div_S(E)$ and define
\begin{equation}\label{eq:JE}
 J_E=\{f\in\AS:\coef{t^\beta}f=0\text{ for all }\beta\in D_E\}.
\end{equation}

\begin{theorem}[Minimal coefficient-retaining quotient]
\label{thm:minquotient}
The set $J_E$ is an ideal, and
\begin{equation}\label{eq:QE}
 Q_E=\AS/J_E
 \cong\bigoplus_{\beta\in D_E}R\,t^\beta
\end{equation}
as $R$-modules. Multiplication is ordinary convolution followed by
restriction to $D_E$. This quotient is the coarsest algebra quotient that
retains the coefficient maps at every target in $E$: if an ideal $J$ has
\[
 \coef{t^\alpha}f=0\quad\text{for every }f\in J\text{ and }\alpha\in E,
\]
then $J\subseteq J_E$. Thus every algebra quotient through which those
coefficient maps descend surjects onto $Q_E$.

The image of $I_1$ in $Q_E$ is nilpotent of exponent at most $H_E+1$, and
the image of $I_{H_E+1}$ is zero. The Euler derivation $\nabla$ descends
to $Q_E$.
\end{theorem}
\begin{proof}
If $\beta\notin D_E$ and $\gamma\in M$, then
$\beta+\gamma\notin D_E$: otherwise adding its complement to a target
would also make $\beta$ a divisor of that target. The complement of $D_E$
is therefore closed under addition of $M$. It follows that $J_E$ is an
ideal. Restriction to $D_E$ identifies every coset with a unique finite
sum of the displayed basis monomials, proving \eqref{eq:QE} and the
multiplication formula. Associativity is inherited from the quotient;
discarding exponents outside $D_E$ cannot remove a contribution that later
returns to a retained exponent.

For minimality, take $f\in J$ and any $\beta\in D_E$. Choose
$\alpha\in E$ and $\gamma\in M$ with $\beta+\gamma=\alpha$.
Since $J$ is an ideal, $t^\gamma f\in J$. Its coefficient at $\alpha$
is exactly the coefficient of $f$ at $\beta$, and must be zero. Therefore
$f\in J_E$. This proves the universal quotient property for all such
ideals, not only monomial ideals.

Every exponent in $D_E$ has depth at most $H_E$ by
Lemma~\ref{lem:divisors}. Hence $I_{H_E+1}$ maps to zero, and
$I_1^{H_E+1}\subseteq I_{H_E+1}$ proves nilpotence.
Finally, $\nabla$ preserves each support exponent. Thus it preserves
$J_E$ and induces a $K$-derivation on the quotient.
\end{proof}

The qualifier ``finite'' refers to rank over $R=K[L]$. Polynomial blocks
can have arbitrary degree, so $Q_E$ need not be finite-dimensional over
$K$. This is exactly the useful distinction for power--logarithmic
computation: finitely many exponent blocks, each represented by a finite
polynomial.

\begin{corollary}[Finite residual certificate]\label{cor:residual}
For a system \eqref{eq:A}, let $\overline\Psi$ be its induced system in
$Q_E$. It has a unique solution in the image of $I_1^d$.
Any candidate $\overline U$ in that ideal satisfying
\[
 \overline\Psi(\overline U)=0
\]
has the exact target coefficients of $U_*$. Only input nonlinear degrees
$|\nu|\le H_E-1$ can contribute when $H_E\ge1$.
The Newton iterate with $2^{n+1}-1>H_E$ automatically supplies such a
candidate. For reversion, the Euler formula \eqref{eq:EulerLB} truncates
exactly after $n=H_E$ in this quotient.
\end{corollary}
\begin{proof}
The image of the original solution is a solution. The Lipschitz estimate
raises support depth by one, and after $H_E+1$ increases the error is
zero in $Q_E$; hence a solution in the augmentation ideal is unique.
A nonlinear term of degree $|\nu|$ has support depth at least
$|\nu|+1$, so terms with $|\nu|\ge H_E$ vanish. The Newton and Lagrange
claims follow from their support-depth bounds.
\end{proof}

This provides an independently checkable certificate: the checker need not
reproduce the Newton trajectory. It can validate the finite dependence
set, the complete relevant input coefficients, and a finite residual
identity in $Q_E$. Such a certificate is only as effective as the supplied
support-completeness proof; the next result makes that proof concrete for
finite generators.

\subsection{A rational certificate for finite supports}

\begin{theorem}[Finite supports admit a positive rational grading]
\label{thm:separator}
Suppose $S=\{s_1,\ldots,s_k\}$ is finite and positive. The subgroup
$H=\langle S\rangle\subseteq\Gamma$ is a finitely generated free abelian
group. There exists an additive map
\[
 w:H\longrightarrow\Q,\qquad w(s_i)\ge1\quad(1\le i\le k).
\]
This map need not preserve the ambient order. For every $\alpha\in M(S)$,
\begin{equation}\label{eq:weightbound}
 \depth_S(\alpha)\le\lfloor w(\alpha)\rfloor.
\end{equation}
Given integer coordinates for the $s_i$ and a rational vector representing
$w$, checking the $k$ inequalities supplies a finite certificate of this
word-length bound.
\end{theorem}
\begin{proof}
Ordered abelian groups are torsion free, so their finitely generated
subgroups are free abelian. Choose a basis of $H$ and regard $s_i$ as
integer vectors in $\R^r$. Their convex hull does not contain zero.
Indeed, a convex representation of zero by rational vectors would admit
one with rational coefficients: reduce to an affinely independent subset
of minimal size and solve its rational linear system. Clearing denominators
would give a nontrivial relation
$\sum_i n_is_i=0$ with all $n_i\ge0$, impossible because every $s_i>0$ in
$\Gamma$.

A compact convex polytope not containing zero is strictly separated from
zero by a real linear functional. One direct construction uses its
nearest point $p$ to zero: for every point $x$ of the polytope,
$p\cdot x\ge\|p\|^2>0$. Approximate the finitely many coefficients of
this functional by rationals closely enough to preserve its strict
positivity at the finitely many $s_i$. Then multiply by a positive rational
constant so that all values are at least one. This gives $w$.
For a word of length $n$ summing to $\alpha$, additivity gives
$w(\alpha)\ge n$, proving \eqref{eq:weightbound}.
\end{proof}

Let the integer coordinate matrix have columns $s_i$. For a target
$\alpha\in H$, all its generator-count representations occur among
\begin{equation}\label{eq:enumeration}
 m\in\N^k,\qquad \sum_i m_i\le\lfloor w(\alpha)\rfloor,
 \qquad \sum_i m_is_i=\alpha.
\end{equation}
This is a finite search. It gives the exact depth by maximizing $\sum_i m_i$.
The ordered words need not all be listed: for each count vector, its
subcount vectors generate all corresponding divisor exponents. Their union
over the target representations is exactly $D_E$.

\subsection{A practical certificate format}
For finite-support reversion, a support certificate can contain an integer
basis for $H$, the generator columns, the rational row $w$, the target
coordinates, and the complete count-vector lists in
\eqref{eq:enumeration}. An exact checker verifies the matrix identities,
the inequalities $w(s_i)\ge1$, and exhaustive enumeration within the finite
budget. It then constructs $D_E$, its projected multiplication table, and
the induced Euler operator from the character values and polynomial
coefficient differentiation.

A solution certificate is a vector of polynomial blocks indexed by $D_E$.
The checker verifies that its exponent-zero block is zero and its residual
in the finite quotient is zero. Corollary~\ref{cor:residual} then certifies
the requested coefficients. Alternatively, the checker may evaluate the
finite Euler formula directly and compare it with a Newton output.

This design cleanly separates trusted finite arithmetic from support
reasoning. It does not assert polynomial total running time. The number
of exponent blocks, the number of support representations, and the bit
sizes of coefficients can be large. The logarithmic claim concerns the
number of Newton rounds in \eqref{eq:rounds}, not the cost of a round or the
size of the final answer.

\begin{example}[A resonant finite quotient]\label{ex:finitequotient}
In lexicographic $\Z^2$, take
\[
 S=\{(0,1),(1,-2),(1,0)\},\qquad \alpha=(2,1).
\]
All generators are positive, although one second coordinate is negative.
The functional $w(a,b)=3a+b$ takes values $1,1,3$ on them and gives
$w(\alpha)=7$. The three count vectors representing $\alpha$ are
\[
 (1,0,2),\quad(3,1,1),\quad(5,2,0).
\]
Their lengths are $3,5,7$, so $\depth_S(\alpha)=7$. The monoid is
$\{(a,b):a\ge0,\ b\ge-2a\}$. Its target-divisor set is
\[
 D_{\{\alpha\}}=
 \{(a,b):a=0,1,2,\ -2a\le b\le5-2a\},
\]
which has $18$ elements. Thus the exact coefficient problem factors through
an $18$-dimensional free module over $K[L]$, even though the ambient
valuation cutoff below $(2,1)$ contains infinitely many powers of $(0,1)$.
Three Newton updates, reaching depth $15>7$, suffice uniformly in this
quotient. The count-vector and quotient-rank assertions are checked by
\pathcode{verify.py} using exact integer arithmetic.
\end{example}

\section{An infinite rank-two support with quadratic depth}
\label{sec:infinite}

The positive rational grading of Theorem~\ref{thm:separator} cannot in
general be chosen on an infinite positive well-ordered support. A concrete
family also shows how an exact nonlinear depth function replaces it.

Let $\Gamma=\Z\eta\oplus\Z\delta$ be lexicographically ordered, with
$\eta=(1,0)$ and $\delta=(0,1)$. Define
\begin{equation}\label{eq:infiniteS}
 S=\{\delta\}\cup\{s_n:n\ge1\},\qquad
 s_n=n\eta-n^2\delta.
\end{equation}
The set is positive and well ordered: $\delta$ is first, and the $s_n$ are
strictly increasing by their first coordinates. It has order type
$1+\omega$.

\begin{theorem}[Exact monoid and nonlinear depth]\label{thm:quadraticdepth}
For the support \eqref{eq:infiniteS},
\begin{equation}\label{eq:infiniteM}
 M(S)=\{k\delta:k\ge0\}
 \ \cup\ \{n\eta+k\delta:n\ge1,\ k\ge-n^2\}.
\end{equation}
Its depth is
\begin{equation}\label{eq:quadraticdepth}
 \depth_S(k\delta)=k\quad(k\ge0),\qquad
 \boxed{\ \depth_S(n\eta+k\delta)=n^2+k+1\ }
 \quad(n\ge1,\ k\ge-n^2).
\end{equation}
There is no additive map $w:\Gamma\to\R$ that is strictly positive on
all of $S$.
\end{theorem}
\begin{proof}
Consider a word with $a$ copies of $\delta$ and $r\ge1$ remaining letters
$s_{n_1},\ldots,s_{n_r}$. Its coordinates are
\[
 n=\sum_{j=1}^r n_j,\qquad
 k=a-\sum_{j=1}^r n_j^2.
\]
Since $\sum n_j^2\le n^2$, it has $k\ge-n^2$. Conversely, every point
with $n\ge1$ and $k\ge-n^2$ is represented by one $s_n$ and $k+n^2$
copies of $\delta$. Words with first coordinate zero consist only of
$\delta$'s. This proves \eqref{eq:infiniteM}.

For the displayed word, its length is
\[
 a+r=k+\sum_{j=1}^r n_j^2+r.
\]
For positive integers with sum $n$,
\[
 n^2-\sum_{j=1}^r n_j^2
 =2\sum_{i<j}n_in_j\ge r-1.
\]
For $r=1$ the inequality is equality; for $r\ge2$ it follows, for example,
from $n_in_j\ge1$ and $r(r-1)\ge r-1$.
Thus every representation length is at most $n^2+k+1$. The one-letter
choice $s_n$ followed by $k+n^2$ copies of $\delta$ attains the bound.
For $n=0$ the depth is plainly $k$. This proves \eqref{eq:quadraticdepth}.

Finally, if an additive $w$ were positive on $S$, put
$c=w(\delta)>0$ and $d=w(\eta)\in\R$. Then
\[
 w(s_n)=nd-n^2c<0
\]
for all sufficiently large $n$, a contradiction.
\end{proof}

The obstruction is not simply that $S$ is infinite, or that its exponent
group has high rank. It is the impossibility of assigning finite real
weights compatible with this entire infinite family. The depth function
is still finite at every exponent and is explicitly computable.


\begin{corollary}[A target-adaptive sharp rational certificate]
\label{cor:adaptive}
For the quadratic support, fix a target $n\eta+k\delta$ with $n\ge1$.
All its word representations use only
$\{\delta,s_1,\ldots,s_n\}$. The rational functional
\[
 w_n(\delta)=1,\qquad w_n(\eta)=n+\frac1n
\]
is at least one on these generators and gives the exact target depth:
\[
 w_n(n\eta+k\delta)=n^2+k+1.
\]
\end{corollary}
\begin{proof}
First coordinates are nonnegative and add, so a word representing the
chosen target cannot use $s_j$ with $j>n$. For $1\le j\le n$,
\[
 w_n(s_j)-1=(n-j)\left(j-\frac1n\right)\ge0.
\]
The asserted target value follows by substitution. It agrees with the
attained depth in Theorem~\ref{thm:quadraticdepth}.
\end{proof}

Thus a single positive real grading is impossible globally, while a sharp
rational certificate exists for every finite target individually. The
functional is allowed to depend on the target; suppressing that dependence
would change the theorem.

\subsection{Arbitrarily fast depth profiles in rank two}
The quadratic example belongs to a larger family.

\begin{theorem}[Superadditive depth-profile realization]
\label{thm:profiles}
Let $(b_n)_{n\ge1}$ be positive integers satisfying
\[
 b_{m+n}\ge b_m+b_n\qquad(m,n\ge1).
\]
In lexicographic $\Z\eta\oplus\Z\delta$, put
\[
 S_b=\{\delta\}\cup\{n\eta-(b_n-1)\delta:n\ge1\}.
\]
This is a positive well-ordered support. Its monoid and depth are
\begin{align*}
 M(S_b)&=\{k\delta:k\ge0\}\ \cup\
 \{n\eta+k\delta:n\ge1,\ k\ge1-b_n\},\\
 \depth_{S_b}(k\delta)&=k,\\
 \depth_{S_b}(n\eta+k\delta)&=k+b_n
       \quad(n\ge1,\ k\ge1-b_n).
\end{align*}
A real-valued additive functional strictly positive on all of $S_b$ exists
if and only if $\sup_n b_n/n<\infty$.
\end{theorem}
\begin{proof}
Positivity and well ordering follow from the first coordinates exactly as
for \eqref{eq:infiniteS}. A word with $a$ copies of $\delta$ and $r\ge1$
other letters of indices $n_1,\ldots,n_r$ has
\[
 n=\sum_jn_j,\qquad k=a-\sum_j(b_{n_j}-1),\qquad
 \text{length}=k+\sum_j b_{n_j}.
\]
Repeated superadditivity gives $\sum_j b_{n_j}\le b_n$.
It also gives $k=a+r-\sum_j b_{n_j}\ge1-b_n$.
Conversely, a single letter of index $n$ and $k+b_n-1$ copies of
$\delta$ realize every permitted $(n,k)$ and attain length $k+b_n$.
Words with first coordinate zero contain only $\delta$'s.
This proves the monoid and depth formulas.

If a functional is positive on $S_b$, write its values on $\delta,\eta$
as $c>0,d\in\R$. Positivity requires
$b_n-1<nd/c$, so $b_n/n<d/c+1/n$ is bounded above.
Conversely, if $M=\sup_n b_n/n$ is finite, the functional with
$w(\delta)=1$ and $w(\eta)=M+1$ satisfies
\[
 w\bigl(n\eta-(b_n-1)\delta\bigr)=n(M+1)-b_n+1\ge n+1>0.
\]
\end{proof}

\begin{corollary}[No rank-two growth ceiling]\label{cor:arbitrarygrowth}
Given any function $g:\N_{\ge1}\to\N_{\ge1}$, there is a positive
well-ordered support in lexicographic $\Z^2$ such that
\[
 \depth_S(n\eta)\ge g(n)\quad\text{for every }n\ge1.
\]
\end{corollary}
\begin{proof}
Set $c_n=\max_{1\le j\le n}g(j)$ and $b_n=nc_n$.
Because $c_n$ is nondecreasing,
\[
 b_{m+n}=(m+n)c_{m+n}\ge mc_m+nc_n=b_m+b_n.
\]
Theorem~\ref{thm:profiles} gives depth $b_n\ge g(n)$ at $n\eta$.
\end{proof}

For example, $b_n=n!$ is superadditive: for $m,n\ge1$,
$(m+n)!\ge2m!n!\ge m!+n!$. It gives the exact depth
$\depth_{S_b}(n\eta)=n!$. The quadratic support corresponds to
$b_n=n^2+1$, which is also superadditive. These examples show that finite
coefficient dependence does not imply a rank-dependent polynomial bound
on the required nonlinear order.

\begin{proposition}[The maximum depth can be noncancelling]
\label{prop:noncancelling}
For any positive well-ordered $S$, put
$\sigma=\sum_{s\in S}t^s$ over the coefficient field $\Q$.
The small-coefficient equation
\[
 V=\frac{\sigma}{1-V}
\]
has solution
\[
 V=\sum_{m\ge1}\frac1m\binom{2m-2}{m-1}\sigma^m.
\]
For every nonzero $\alpha\in M(S)$, the term with
$m=\depth_S(\alpha)$ makes a strictly positive contribution to
$\coef{t^\alpha}V$.
\end{proposition}
\begin{proof}
The equation is equivalent to $V^2-V+\sigma=0$; its zero-constant branch
gives the displayed Catalan series. Positive-word summability justifies
specializing an ordinary formal variable to $\sigma$.
The coefficient of $t^\alpha$ in $\sigma^m$ counts ordered support words
of length $m$ with that sum. It is a nonnegative integer, finite by the
word lemma, and is positive when $m=\depth_S(\alpha)$. The multiplying
Catalan coefficient is positive. Thus the maximum-depth contribution is
realized without cancellation.
\end{proof}

\subsection{A genuine infinite-support inverse}
Choose $h=\delta$ and any character with $\chi(\delta)=1$, for example
$\chi(n\eta+k\delta)=n\omega+k$ over $K=\Q(\omega)$. Let
\begin{equation}\label{eq:infiniteB}
 b=t^\delta+\sum_{n\ge1}t^{n\eta-n^2\delta},\qquad
 B=Xb,
 \qquad F=X+B.
\end{equation}
Then $v(b)=\delta>0$, so Theorem~\ref{thm:reversion} applies. The inverse
has normalized support in \eqref{eq:infiniteM}. For any $n\ge1$ and
$k\ge-n^2$, its coefficient at $n\eta+k\delta$ is determined by the
Lagrange orders
\[
 1\le m\le n^2+k+1,
\]
and by Newton stage
\begin{equation}\label{eq:infiniteRounds}
 j=\left\lfloor\log_2(n^2+k+2)\right\rfloor.
\end{equation}
These are guaranteed support bounds, not assertions that every possible
coefficient is nonzero or that cancellation cannot improve the round count
for a particular coefficient.

There is also a direct effective enumeration. A target with first
coordinate $n$ can use only $s_1,\ldots,s_n$. Enumerate integer partitions
of $n$ to determine their multiplicities; the required number of
$\delta$'s is then $a=k+\sum n_i^2$, which must be nonnegative.
This finite calculation gives the exact depth and the relevant divisor
set, without searching all support terms below a valuation cut. The
verification program independently performs this partition enumeration
for $1\le n\le10$ and $-n^2\le k\le10$.

\begin{remark}[What this example rules out]
One cannot replace the arbitrary-group theory by an argument that always
first assigns a positive real additive weight to every input exponent.
Such a weight exists for finite supports but not for \eqref{eq:infiniteS}.
The word-depth proof and the finite target quotient remain valid, so this
is an obstruction to a proposed proof and precision representation, not
to inversion itself.
\end{remark}

\section{Exact verification and reproducibility}
\label{sec:verification}

\subsection{What was executed}
The accompanying \pathcode{verify.py} uses Python's exact rational
arithmetic and polynomial arithmetic over $\Q[L]$ in SymPy. It performs no
floating-point tests and makes no network requests. The recorded run used
Python 3.13.5 and SymPy 1.14.0. It writes
\pathcode{verification_results.json}, including the tested ranges,
coefficients, support lists, and assertions.

Four groups of checks were executed successfully. They concern the scalar
sharpness example, a resonant power--logarithmic inverse, the finite
quotient in Example~\ref{ex:finitequotient}, and the quadratic-depth formula
in Theorems~\ref{thm:quadraticdepth} and~\ref{thm:profiles}. These are finite audits of signs,
indexing, normalization, resonance handling, and support enumeration.
The quantified theorems depend on the proofs, not on extending a finite
pattern by assumption.

\subsection{Scalar sharpness through degree 130}
The program represents ordinary power series through degree $130$ by
arrays of rational coefficients. It checks the exact root residual, then
performs Newton updates for the rational equation
$u+z/(1+u)=0$. The first nonzero errors are:
\begin{center}
\begin{tabular}{@{}rrrr@{}}
\toprule
Newton index $n$ & Error degree $r_n$ & Leading coefficient & $(2r_n-1)$\\
\midrule
0 & 1 & 1 & 1\\
1 & 3 & 1 & 5\\
2 & 7 & 1 & 13\\
3 & 15 & 1 & 29\\
4 & 31 & 1 & 61\\
5 & 63 & 1 & 125\\
6 & 127 & 1 & 253\\
\bottomrule
\end{tabular}
\end{center}
The last column is the multiple of $h$ in the displacement-error valuation
for \eqref{eq:ranktwoF}. Every entry agrees with the proven formulas.

\subsection{A resonant power--logarithmic example}
For an independent block calculation, take lexicographic $\Gamma=\Z^2$,
$h=(0,1)$, and $\chi(a,b)=b$. Let
\[
 s_1=(0,2),\quad s_2=(1,-3),\quad s_3=s_1+s_2=(1,-1),
 \qquad q_i=t^{s_i}\ (i=1,2),
\]
and set
\begin{equation}\label{eq:verifiedb}
 b=q_1(L+1)+q_2(L^2-2)+q_1q_2L.
\end{equation}
The third input block is a primitive coefficient at the same exponent as
the product of the first two generators. It therefore tests a genuine
support resonance, rather than a freely indexed three-variable model.
The normalized feedback equation is
\begin{align*}
 A_B(U)={}&q_1(1+U)^{-1}\bigl(L+\log(1+U)+1\bigr)\\
 &+q_2(1+U)^4\bigl((L+\log(1+U))^2-2\bigr)\\
 &+q_1q_2(1+U)^2\bigl(L+\log(1+U)\bigr).
\end{align*}
The program computes all $27$ nonconstant blocks of total $q_1,q_2$
degree at most six. It compares the Euler--Lagrange formula, six Picard
updates, and two Newton updates. All three results agree exactly. Both
composition residuals vanish through the cutoff. In particular, it checks
not only $F\circ G=X$ but also $G\circ F=X$, whose normalized residual is
\[
 b+(1+b)T_B(U).
\]
\Needspace{12\baselineskip}
Some resulting blocks are:
\begin{center}
\begin{tabular}{@{}lp{0.74\textwidth}@{}}
\toprule
Monomial & Coefficient in $U$\\
\midrule
$q_1$ & $-L-1$\\[0.2em]
$q_2$ & $2-L^2$\\[0.2em]
$q_1q_2$ & $3L^3+6L^2-5L-8$\\[0.2em]
$q_1^2$ & $-L^2-L$\\[0.2em]
$q_2^2$ & $4L^4+2L^3-16L^2-4L+16$\\[0.2em]
$q_1^2q_2$ & $-L^4-8L^3-13L^2+4L+10$\\
\bottomrule
\end{tabular}
\end{center}
No finite total-degree cutoff is claimed to represent the full valuation
cut below an arbitrary lexicographic exponent. It is the finite quotient
appropriate to this particular block audit.

\subsection{Support and depth checks}
For Example~\ref{ex:finitequotient}, exhaustive enumeration under the
certified weight budget $7$ gives exactly the three stated count vectors,
maximum length $7$, and $18$ divisor exponents. A separate downward-closure
check confirms that projected multiplication cannot discard a future
contribution to a target.

For \eqref{eq:infiniteS}, the program enumerates the integer partitions of
$n$ and computes each feasible representation length independently of the
closed formula. It checks $495$ pairs with
$1\le n\le10$ and $-n^2\le k\le10$. All maximum lengths equal
$n^2+k+1$. A further exact partition audit tests the factorial profile
$b_n=n!$ at $49$ selected boundary and interior targets for $1\le n\le9$.

\subsection{Rebuilding the package}
From the package directory, run
\begin{verbatim}
python verify.py
pdflatex -interaction=nonstopmode -halt-on-error article.tex
pdflatex -interaction=nonstopmode -halt-on-error article.tex
pdflatex -interaction=nonstopmode -halt-on-error article.tex
\end{verbatim}
The supplied build script performs the same steps. The TeX source contains
its bibliography and tables; it does not require the repository checkout,
external images, or a bibliography database. The verification output is a
separate audit record rather than a hidden source of unstated hypotheses.

\section{A statement-level formalization plan}
\label{sec:formalization}

The inspected Lean module
\pathcode{Analysis/Transseries/Lean/Transseries/TransseriesWellBased.lean}
contains the binary support bridge: the named multiplicative declarations
\pathcode{neumann_isPWO} and
\pathcode{neumann_finite_factorizations}, their generated additive
counterparts \pathcode{neumann_add_isPWO} and
\pathcode{neumann_finite_decompositions}, Dickson-type declarations, and
order-dual wrappers. Its namespace remains \pathcode{Fabius}
\cite{RepoLean}. These are relevant foundations, not a formal proof of
this article's inverse or Newton theorem.

\Needspace{22\baselineskip}
A realistic extension proceeds through mathematical interfaces rather than
marking an entire chapter as formalized because it mentions Hahn series.
\begin{center}
\begin{tabular}{@{}p{0.23\textwidth}p{0.70\textwidth}@{}}
\toprule
Layer & Statements to formalize\\
\midrule
All-word support & Higman embedding for finite words; positive-sum
strictness; well-ordering of $M(S)$; finite all-length fibers.\\[0.4em]
Depth filtration & Maximum representation length; support ideals $I_r$;
separatedness; completeness and strong summation of rising-depth families.\\[0.4em]
Implicit systems & Summable evaluation; the one-factor and two-factor
remainder estimates; unipotent matrix inverse; Newton depth $2r+1$.\\[0.4em]
Differential blocks & Character data $(h,\chi)$; strong derivative;
iterated block identity; induced support-preserving Euler derivation.\\[0.4em]
Substitution and reversion & Strong Taylor map; generator identities;
associativity; justified parameter specialization; normalized Euler formula.\\[0.4em]
Finite certificates & Divisor closure; largest coefficient-annihilating
ideal; finite quotient multiplication; nilpotence; residual checker.\\[0.4em]
Sharp boundaries & Exact scalar Newton error identity; noncofinal rank-two
counterexample; uniform valuation criterion; quadratic-depth family.\\
\bottomrule
\end{tabular}
\end{center}

The finite quotient theorem suggests a particularly manageable first
implementation target. A checker for finite exponent lists, rational
characters, polynomial blocks, and a supplied rational weight certificate
can be separated from the more difficult abstract Hahn infrastructure.
Its soundness theorem would explicitly invoke the completed support
lemmas, rather than infer them from finite examples.

No new Lean module was compiled for this package. The exact Python checks
are reproducible executable tests, and the mathematical arguments are
conventional proofs. The distinction is important both for the current
status report and for deciding which formalization dependencies remain.

\section{Further research questions}
\label{sec:questions}

The questions below are proposed continuations with explicit mathematical
targets. Their appearance here does not assert that every formulation is
an established open problem in the whole literature.

\begin{question}[Variable-shift derivations and genuine multi-scale calculus]
Replace \eqref{eq:D} by a strong derivation for which differentiating a
monomial can produce several exponent shifts depending on that monomial.
Which well-founded condition on the resulting support-dependency graph
replaces the common positive monoid? Seek a substitution theorem with
finite anchored-dependency fibers and a computable analogue of
$\depth_S$, without requiring all shifts to lie in one Archimedean class.
\end{question}
The current proof uses a fixed shift to cancel the $-h$ normalization in
every Taylor order. A variable-shift extension needs a replacement for that
cancellation, not merely a change of notation.

\begin{question}[When does ordinary sequential iteration suffice?]
For the small-coefficient systems here, an ordinary sequence indexed by
$\N$ reaches every coefficient; transfinite Newton stages are unnecessary.
Characterize strong nonlinear operators outside this class for which the
same statement remains true. Can one isolate an intrinsic finite
coefficient-dependency condition equivalent to stabilization after ordinary
sequential iteration, and distinguish it from the existence of an abstract
fixed point obtained by a different construction?
\end{question}

\begin{question}[Cancellation-sensitive precision]
The depth $\depth_S$ is a support-only invariant and ignores coefficients.
Find efficiently checkable hypotheses under which the coefficient at a
maximum-depth target is nonzero, and a certificate format for cancellations
that reduce the necessary Newton depth. The scalar sharpness theorem shows
that no unconditional improvement is possible, but structured families
may admit much smaller certified dependence sets.
\end{question}

\begin{question}[Effective infinite supports]
For which computable descriptions of infinite well-ordered supports can
one decide membership in $M(S)$, compute $\depth_S(\alpha)$, and enumerate
$\Div_S(E)$? Presburger descriptions, finite-state descriptions, and
polynomially parameterized lexicographic supports are natural test classes.
The quadratic family gives a nontrivial solvable example; existential
finiteness alone is not an algorithm.
\end{question}

\begin{question}[Compressed infinite cuts]
When Theorem~\ref{thm:cut} rules out a finite monomial list, determine when
a cutoff nevertheless admits a finite algebraic, rational, or differential
block representation that is stable under substitution. A useful objective
is a certificate calculus whose objects combine finite divisor quotients
with exact infinite lower-rank blocks, so that crossing a noncofinal cut
does not force term-by-term enumeration.
\end{question}

\begin{question}[Bit complexity of minimal coefficient quotients]
Give output-sensitive complexity bounds for constructing $D_E$, multiplying
in $Q_E$, and checking residual certificates. Determine when resonance
compression saves work compared with a free multi-index expansion. The
logarithmic number of Newton rounds is established here; polynomial total
complexity is not.
\end{question}

\begin{question}[Analytic realization and differential characters]
Identify coefficient fields, characters, support classes, and evaluation
maps for which the formal substitution and inverse agree with operations
on actual Hardy-field or sectorial analytic germs. In a non-Archimedean
support hierarchy, an injective character into an abstract coefficient
field need not be an analytic exponent assignment. A successful realization
theorem must address both the ordered monomial hierarchy and the coefficient
field, not only the formal support estimates.
\end{question}

\begin{question}[Divided powers and positive characteristic]
The word-support and finite-quotient arguments do not intrinsically require
factorials. Can the reversion theorem be expressed using Hasse--Schmidt
operators or divided-power structures so as to isolate exactly what
survives in positive characteristic? In particular, distinguish the
support-theoretic implicit theorem from the characteristic-zero Euler and
Lagrange formulas.
\end{question}

\begin{question}[An independently checked finite certificate engine]
Implement the quotient certificate of Section~\ref{sec:certificates} in a
proof assistant, beginning with integer-coordinate exponents, rational
characters, and polynomial logarithmic blocks. The first milestone should
be a kernel-checked implication from a complete finite support certificate
and a zero residual to exact target coefficients, including resonances.
A subsequent milestone is certificate synthesis rather than merely
certificate checking.
\end{question}

\section{Conclusion}
The repository's non-Archimedean reversion question has a positive answer
in the stated fixed-shift Hahn power--logarithmic setting: cofinal valuation
improvement is not needed for existence, the ordinary integer Lagrange
index is sufficient, and the inverse remains inside the normalized input
support monoid. The correct coefficientwise finiteness invariant is
maximum positive-word length.

That invariant produces more than an existence proof. It yields the sharp
Newton law $r\mapsto2r+1$, an exact uniform criterion for valuation
convergence, and a minimal finite quotient for any finite coefficient
request. The quadratic-depth example demonstrates that no global positive
real additive grading can always replace this combinatorics, even in a
rank-two group. The superadditive-profile construction allows arbitrary
depth growth in that same rank. The result is a support and precision calculus that
separates valid formal computation from stronger topological or analytic
claims and supplies a concrete next target for verified symbolic
implementation.

\appendix
\section{The coefficient Lagrange identity over a general \texorpdfstring{$\Q$}{Q}-algebra}
\label{app:lagrange}

For clarity, we record a proof that does not assume the constant coefficient
of the relevant series is invertible after specialization.

\begin{lemma}[Formal coefficient inversion]\label{lem:ordinaryLB}
Let $A$ be any commutative $\Q$-algebra and $\phi(y)\in A[[y]]$.
There is a unique $c(z)\in zA[[z]]$ satisfying $c=z\phi(c)$, and
\[
 [z^n]c(z)=\frac1n[y^{n-1}]\phi(y)^n\qquad(n\ge1).
\]
\end{lemma}
\begin{proof}
Successive comparison of coefficients gives existence and uniqueness:
the coefficient at $z^n$ on the right depends only on coefficients of
$c$ of degrees less than $n$. Moreover, each coefficient is a polynomial
with rational coefficients in finitely many coefficients of $\phi$.

First work in the universal polynomial ring
$A_0=\Q[a_0,a_1,a_2,\ldots]$, with
$\phi(y)=\sum_{j\ge0}a_jy^j$, and localize by inverting $a_0$.
Then $z=y/\phi(y)$ is a formal change of parameter with invertible linear
coefficient. Formal residues obey change of variable:
\[
 \Res_z f(z)\,dz=\Res_y f(z(y))z'(y)\,dy.
\]
For Laurent monomials other than $z^{-1}$ this follows because the
substituted expression is a derivative and has zero residue; for
$z^{-1}$ the logarithmic derivative has residue one. Linearity gives the
identity for a Laurent series with finite principal part, as used below.

Using $c(z(y))=y$, we obtain
\begin{align*}
 [z^n]c(z)
 &=\Res_z c(z)z^{-n-1}\,dz\\
 &=\Res_y y^{-n}\bigl(\phi(y)^n-y\phi(y)^{n-1}\phi'(y)\bigr)\,dy\\
 &=[y^{n-1}]\phi(y)^n-[y^{n-2}]\phi(y)^{n-1}\phi'(y)\\
 &=\frac1n[y^{n-1}]\phi(y)^n.
\end{align*}
The last equality follows by differentiating $\phi^n$; it also holds at
$n=1$ with the negative-index coefficient interpreted as zero.
Both sides of the resulting coefficient identity lie in $A_0$, and
$A_0$ injects into its localization. Hence the identity already holds in
$A_0$. Specialize its finitely many relevant indeterminates to any
coefficients in $A$ to obtain the asserted formula without any invertibility
condition on $\phi(0)$. Replacing $\phi$ by $-\phi$ supplies the sign
$(-1)^n$ used in \eqref{eq:LB}.
\end{proof}

\Needspace{28\baselineskip}
\section{Source and verification status ledger}
\label{app:ledger}

\begin{longtable}{@{}p{0.31\textwidth}p{0.62\textwidth}@{}}
\toprule
Item & Status and boundary\\
\midrule
\endfirsthead
\toprule
Item & Status and boundary\\
\midrule
\endhead
Repository snapshot & Pinned to commit
\pathcode{3f2d57344dc6f0b3dfa376f9490bfe13d1beb81e}.
The recent reversion question, README, and support Lean module were
inspected; no exhaustive repository audit is claimed.\\[0.5em]
Classical support background & Higman word embeddings, Hahn finite-fiber
summation, and ordinary coefficient Lagrange inversion. Proofs of the
specific versions used here are included.\\[0.5em]
Conventional proofs & Theorems on support depth, implicit systems,
fixed-shift substitution and reversion, Newton precision, valuation and
cut criteria, minimal finite quotients, and the quadratic-depth family.\\[0.5em]
Executed exact checks & Scalar Newton through ordinary degree $130$;
resonant block inversion through total degree $6$; the $18$-block finite
quotient; $495$ quadratic-profile and $49$ factorial-profile depth comparisons.\\[0.5em]
Not executed or not claimed & No new Lean compilation, no universal
analytic realization, no Borel or resurgence theorem, no complete
literature-priority determination, and no polynomial-time complexity
claim.\\[0.5em]
Reproducible materials & Standalone \pathcode{article.tex}, compiled PDF,
\pathcode{verify.py}, exact JSON results, build script, requirements,
README, and provenance record.\\
\bottomrule
\end{longtable}

\begingroup
\small
\begin{thebibliography}{99}
\bibitem{RepoReversion}
ProveIt research corpus,
\emph{Support-Controlled Reversion and the Exact One-Exponential
Substitution Group}, 29 September 2026.
In particular, the further-research question on non-Archimedean exponent
groups and the question on proof-producing support inference.
\href{https://github.com/VladimirReshetnikov/ProveIt/blob/3f2d57344dc6f0b3dfa376f9490bfe13d1beb81e/Analysis/Transseries/docs/series-and-transseries/Support_Controlled_Reversion_One_Exponential/reversion_and_one_exponential.tex}{Pinned TeX source, commit \texttt{3f2d57344dc6}}.

\bibitem{RepoReadme}
ProveIt contributors, \emph{Transseries}, repository README and
series-and-transseries documentation README, snapshot of 29 September 2026.\\
\href{https://github.com/VladimirReshetnikov/ProveIt/blob/3f2d57344dc6f0b3dfa376f9490bfe13d1beb81e/Analysis/Transseries/README.md}{Project README};
\href{https://github.com/VladimirReshetnikov/ProveIt/blob/3f2d57344dc6f0b3dfa376f9490bfe13d1beb81e/Analysis/Transseries/docs/series-and-transseries/README.md}{documentation README}.

\bibitem{RepoLean}
ProveIt contributors, \emph{TransseriesWellBased.lean},
Lean source, same pinned snapshot.\\
\href{https://github.com/VladimirReshetnikov/ProveIt/blob/3f2d57344dc6f0b3dfa376f9490bfe13d1beb81e/Analysis/Transseries/Lean/Transseries/TransseriesWellBased.lean}{Pinned Lean source}.

\bibitem{Higman}
G.~Higman, Ordering by divisibility in abstract algebras,
\emph{Proceedings of the London Mathematical Society} (3) \textbf{2}
(1952), 326--336.
\href{https://doi.org/10.1112/plms/s3-2.1.326}{doi:10.1112/plms/s3-2.1.326}.

\bibitem{vdH}
J.~van der Hoeven, Operators on generalized power series,
\emph{Illinois Journal of Mathematics} \textbf{45} (2001), 1161--1190.
\href{https://doi.org/10.1215/ijm/1258138061}{doi:10.1215/ijm/1258138061}.
Author's version:\\
\url{https://www.texmacs.org/joris/noeth/noeth.html}.

\bibitem{Edgar}
G.~A.~Edgar, \emph{Transseries: Composition, Recursion, and Convergence},
arXiv:0909.1259.\\
\url{https://arxiv.org/abs/0909.1259}.

\bibitem{BM}
V.~Bagayoko and V.~Mantova,
\emph{Taylor expansions over generalised power series},
arXiv:2509.08473 (2025).\\
\url{https://arxiv.org/abs/2509.08473}.

\bibitem{Gessel}
I.~M.~Gessel, Lagrange inversion,
\emph{Journal of Combinatorial Theory, Series A} \textbf{144} (2016),
212--249.
\href{https://doi.org/10.1016/j.jcta.2016.06.018}{doi:10.1016/j.jcta.2016.06.018}.
\url{https://arxiv.org/abs/1609.05988}.

% ed. (2026-09-29): two bibliography entries for the editorial notes.
\bibitem{EdSurcomplex}
ProveIt research corpus, \emph{Surcomplex Analysis: germs, coherent Hahn
sections, and infinitesimal zero geometry}, September 2026,
Lemma \texttt{e:lem-reversion} (Hahn reversion).
Repository path
\path{Algebra/SurrealNumbers/docs/surcomplex/analysis/article.tex}.
Added by the ProveIt editors, 2026-09-29.

\bibitem{EdSurrealLean}
ProveIt contributors, \emph{Surreal.HahnSeries.Neumann} and
\emph{Surreal.HahnSeries.NeumannWords}, Lean sources
\path{Algebra/SurrealNumbers/Surreal/HahnSeries/Neumann.lean} and
\path{Algebra/SurrealNumbers/Surreal/HahnSeries/NeumannWords.lean}.
Added by the ProveIt editors, 2026-09-29.
\end{thebibliography}
\endgroup
\end{document}
